\section{Introduction}\label{sec:introduction}

A division ring $D$ is \emph{algebraic over its center} $F=Z(D)$ if every
$a\in D$ satisfies a nonzero polynomial in $F[t]$. For a subset $S\subseteq D$,
write $F(S)$ for the division subring generated by $F\cup S$.
The ring is \emph{locally finite} if
\[
 [F(S):F]<\infty\qquad\text{for every finite subset }S\subseteq D.
\]
The Kurosh problem for division rings asks whether algebraicity over the
center implies local finiteness. We give a negative answer.

All rings and subrings in this paper have the same identity. If $k$ is a
central subfield and $x,y$ belong to a ring, $k[x,y]$ denotes the unital
$k$-algebra generated by $x,y$; the generators need not commute.
The free polynomial algebra on two formal generators is denoted
$k\langle x,y\rangle$.

\begin{theorem}\label{thm:main}
There is a countable division ring $D$ of characteristic zero, with
center $F=Z(D)$, and elements $x,y\in D$ such that
\[
 D=F[x,y]=F(x,y),\qquad [D:F]=\infty,
\]
while every element of $D$ is algebraic over $F$.
\end{theorem}

The same pair $\{x,y\}$ therefore witnesses the failure of local finiteness.
The construction does not prescribe the center. It first produces an
algebraic division ring over a countable central subfield, and a separate
finite-descent argument proves infinite dimension over the actual center.
The annihilating polynomial is allowed to depend on the element.

\subsection{Historical context}

Kurosh formulated his question in 1941 in connection with Burnside's
finiteness problem for groups. He asked whether a division ring that is
algebraic and finitely generated as an algebra over a central subfield
must be finite-dimensional over that subfield
\cite[Problem~(K), p.~240]{Kurosh1941}. The subfield in this formulation
need not be the full center, and finite generation means algebra
generation, not merely division generation.

The broader problem for algebraic algebras has a different history.
Golod's construction, based on the Golod--Shafarevich estimates, gives
finitely generated infinite-dimensional nil algebras
\cite[Example~1]{Golod1964}. Their unitizations are algebraic algebras
with nonzero nilpotents, so they do not answer the division-ring question.

There are positive results under additional hypotheses. Jacobson's
bounded-degree theorem gives a polynomial identity, and Kaplansky's
theorem makes algebraic algebras satisfying a polynomial identity
locally finite; see \cite[Theorems~3 and~5]{Jacobson1952}.
Bell, Drensky, and Sharifi later proved that a division ring whose
elements satisfy nonzero polynomial relations of uniformly bounded degree $d$,
with coefficients on the left in a possibly noncentral subfield, has
dimension at most $d^2$ over its center
\cite[Theorem~1.3]{BellDrenskySharifi2013}. Our construction imposes
element-dependent relations without a uniform degree bound.

Amitsur proved that a division algebra finitely generated as an algebra
over an uncountable central field is automatically algebraic
\cite[Corollary~$2'$]{Amitsur1956}. This conclusion is algebraicity,
not local finiteness. Countability plays a different role here: it
allows us to schedule all polynomial expressions in the construction.
Goodearl and Letzter's recent work \cite{GoodearlLetzter2026} treats
nonfinite generation as an algebra for classes of transcendental skew
fields and explicitly distinguishes that question from the algebraic
Kurosh problem. Theorem~\ref{thm:main} concerns algebraic division rings
and their actual centers.

The proof combines standard division-algebra constructions with matrix
and formal-series calculations. Symbol and cyclic algebras provide
division-ring models for the generic block data
\cite[Section~2.5]{GilleSzamuely2006}. Traces organize cyclic word
relations \cite[Section~4]{Procesi1976}, while inner derivations describe
infinitesimal conjugation \cite[Theorem~8]{Jacobson1937Derivations}.
The lifting argument supplies the additional response, compatibility,
and generation statements needed to preserve two independent deformation
directions while imposing successive algebraic relations.

\subsection{The construction}

The main ingredient is Theorem~\ref{thm:lifting}. Its input is a tuple of
formal series in two noncommuting variables over a finite-dimensional
central division algebra. The constant tuple generates the coefficient
algebra, and the two linear coefficients are independent modulo changes
coming from conjugation. The theorem makes any prescribed series
algebraic over a larger countable central field. At the same time it
increases the degree of the division algebra, preserves all previous
polynomial identities, and restores both generation and the two independent
directions. These last requirements allow the construction to continue.

For the prescribed series $w$, the immediate goal is to find a formal
substitution $Z=(Z_1,Z_2)$ for which $w(Z)$ satisfies the minimal
polynomial of the chosen constant output $w(Z(0))$. The input $Z(0)$
is a nonzero generic input, small with respect to an auxiliary valuation,
not the original zero input.
After splitting
the coefficient algebra into matrices, we take the constant input to be
a diagonal sum of independent generic matrix pairs. Varying a diagonal
block gives a first eigenvalue response; adding an arrow to another block
and an arrow back gives a paired response. Sections~\ref{sec:formal-generic}
and~\ref{sec:responses} establish the formal evaluation rules and show that
these responses partition the eigenvalue labels into linked groups.

A transfer theorem puts each first unsolved noncommutative word coefficient
in the same response space. It also turns proportional responses into
functional relations between eigenvalue branches
(Sections~\ref{sec:word-transfer} and~\ref{sec:linkage-replication}).
The relations within each block then separate from a fixed-dimensional
space of relations among blocks. To test the latter, we repeat each block
and attach auxiliary blocks, called hubs. Individual hubs distinguish
repeat positions; a common hub compares them. This matrix test proves
that every leading error satisfies the required cyclic sums of word
coefficients (Sections~\ref{sec:global-defect} and~\ref{sec:cyclic-obstructions}).

An eleven-block arrangement of two colors supplies corrections whose
only obstructions are those cyclic sums. We solve the equations one word
length at a time in Section~\ref{sec:free-correction}. A cyclic division
algebra realizes these matrix calculations over a division ring, and a
Taylor translation restores generation and the two independent directions
(Sections~\ref{sec:division-realization} and~\ref{sec:generation}).
The corrections are solutions of finite affine linear systems; this is
what permits their descent from a splitting field to the division
algebra's center. If $w$ already satisfies its constant term's minimal
polynomial, the spectral correction is unnecessary, but realization
and translation are still needed to increase degree and restore the
iteration hypotheses.

Section~\ref{sec:countable-construction} first explains how the lifting
theorem yields Theorem~\ref{thm:main}. The fields introduced during the
iteration are countable, so every polynomial expression can be scheduled.
A division-ring ultraproduct retains the resulting relations. Its
two-generated subalgebra is already a division ring, and finite word
sets prove that its dimension over the central coefficient field is
infinite. The actual-center argument concludes this deduction. The
remaining sections prove the lifting theorem and all its required
compatibilities.
